\section{Threats to Validity}
The collection supports a qualitative map of research directions rather than estimates of their prevalence. Web-based discovery, English-language emphasis, and the stated source restrictions can omit relevant studies. Performance, accessibility, concurrency, and distributed-system testing received less focused coverage than functional testing and fuzzing. Selected neural predecessors provide context but do not constitute a complete history of learned test generation.

Inspection depth also varies across records. Abstract-level descriptions establish broad mechanisms, whereas detailed context permissions, training data, and evaluation denominators require full-text coding. The present synthesis uses those limits to constrain its claims: it reports mechanisms and analytical relationships without ranking methods by numerical gains. A systematic extension should apply a fixed multi-database protocol, complete variant-level extraction, and independent validation of the coding.

Finally, the framework is an analytical organization of the reviewed work. Its usefulness for reproducible classification has not yet been tested through an inter-rater study. Some systems naturally occupy several categories, particularly when they combine scenario generation, execution, and oracle repair. Recording variants and stages separately can preserve these distinctions better than forcing each paper into a single category.

Preprint revisions and changing model services can alter both method descriptions and reproducibility conditions. Publication bias can also favor successful applications over unsuccessful deployments. The dated source register supports version tracking, but it cannot remove these biases; updated searches and explicit treatment of negative evidence remain necessary for a systematic extension.

\section{Conclusion}
LLM-based test generation spans scenarios, executable suites, oracles, and reusable fuzzing infrastructure. We connected these artifacts to their information sources, generation mechanisms, and evaluation assumptions. The synthesis separates executable construction from oracle justification: coverage, reference consistency, patch discrimination, independent detection, and acceptance of legal behavior supply complementary evidence. Making information access, evidence provenance, and cost explicit supports more interpretable comparisons and guides research on generators that remain useful across repositories, adversaries, and software changes.

\section*{Data Availability}
The source register, search notes, synthesis outline, bibliography, and manuscript sources accompany this draft as local supplementary materials. The register includes stable source links and inspection-depth annotations. No experimental dataset or reproduced performance results are claimed. A public archival release has not been created.
